% ============================================================
%  Преамбюл — пакети, шрифтове, език (български), макроси
%  Проектиран за XeLaTeX / Tectonic.
% ============================================================

% --- Шрифтове (Unicode, с поддръжка на кирилица) ----------
% Вграденият пакет на Tectonic не съдържа кирилица в удебелените
% разрези на LM/TeX Gyre, затова ползваме системни шрифтове на macOS
% с пълна кирилица (Times New Roman / Arial / Courier New).
% При компилиране на друга машина смени Path/имената на шрифтовете.
\usepackage{fontspec}
\defaultfontfeatures{Ligatures=TeX}
\defaultfontfeatures[\rmfamily,\sffamily,\ttfamily]{
  Path           = /System/Library/Fonts/Supplemental/ ,
  Extension      = .ttf ,
  UprightFont    = * ,
  BoldFont       = * Bold ,
  ItalicFont     = * Italic ,
  BoldItalicFont = * Bold Italic ,
}
\setmainfont{Times New Roman}
\setsansfont{Arial}
\setmonofont{Courier New}

% --- Език -------------------------------------------------
% babel (а не polyglossia): под XeLaTeX не изисква всеки шрифт (вкл.
% моноширинния) да декларира OpenType script „cyrl'', което иначе чупи
% компилацията при латински думи в заглавия.
\usepackage[english,bulgarian]{babel}

% --- Геометрия и оформление -------------------------------
\usepackage[a4paper,margin=2.5cm]{geometry}
\usepackage{setspace}
\onehalfspacing
\usepackage{indentfirst}      % отстъп и на първия абзац (българска типография)
\usepackage{microtype}

% --- Математика -------------------------------------------
\usepackage{amsmath,amssymb,amsthm}
\usepackage{mathtools}

% --- Фигури, таблици, графики -----------------------------
\usepackage{graphicx}
\graphicspath{{figures/}}
\usepackage{float}
\usepackage{booktabs}
\usepackage{array}
\usepackage{multirow}
\usepackage{caption}
\usepackage{subcaption}
\captionsetup{font=small,labelfont=bf}

% --- Код / листинги ---------------------------------------
\usepackage{listings}
\usepackage{xcolor}
\definecolor{codegray}{gray}{0.95}
\lstset{
  basicstyle=\ttfamily\small,
  backgroundcolor=\color{codegray},
  breaklines=true,
  frame=single,
  numbers=left,
  numberstyle=\tiny,
  captionpos=b,
  showstringspaces=false,
}

% --- Хипервръзки (последно) -------------------------------
\usepackage[hidelinks]{hyperref}

% --- Удобни макроси ---------------------------------------
\newcommand{\eng}[1]{\foreignlanguage{english}{#1}}   % за англ. термини в скоби
\newcommand{\todo}[1]{\textbf{\color{red}[TODO: #1]}}
